\documentclass[11pt]{article}

\usepackage[margin=1in]{geometry}
\usepackage{amsmath,amssymb,amsthm,mathtools}
\usepackage[T1]{fontenc}
\usepackage[utf8]{inputenc}
\usepackage{booktabs}
\usepackage[colorlinks=true,linkcolor=blue,citecolor=blue,urlcolor=blue]{hyperref}

\newtheorem{theorem}{Theorem}
\newtheorem{lemma}[theorem]{Lemma}
\newtheorem{proposition}[theorem]{Proposition}
\newtheorem{corollary}[theorem]{Corollary}
\newtheorem{remark}[theorem]{Remark}

\DeclareMathOperator*{\argmin}{arg\,min}
\newcommand{\ones}{\mathbf{1}}
\newcommand{\x}{\mathbf{x}}
\newcommand{\y}{\mathbf{y}}
\newcommand{\rbeq}{\bar\rho_{=}}
\newcommand{\rble}{\bar\rho_{\le}}
\newcommand{\rseq}{\rho^\star_{=}}
\newcommand{\rsle}{\rho^\star_{\le}}

\title{Penalty Thresholds for the Sensor Placement QUBO:\\
Equality and Inequality Sensor Budgets}
\author{}
\date{}

\begin{document}
\maketitle

\begin{abstract}
We characterize the penalty parameter required for exactness of the QUBO
reformulation of the water distribution network (WDN) sensor placement problem,
for an equality budget $\sum_i x_i=s$ and for an inequality budget
$\sum_i x_i\le s$ encoded with binary slack variables. For arbitrary vertex costs
$c_i$ and nonnegative pipe weights $w_{ij}$, the penalty
$\rho>\bar\rho_{=}:=\max_i\max\{c_i,\,W_i-c_i\}$ is strictly exact for the
equality budget and $\rho>\bar\rho_{\le}:=\max\{0,\max_i(W_i-c_i)\}$ for the
inequality budget, where $W_i$ is the total weight of the pipes at $i$. Unlike the
balanced minimum-cut threshold, these values do not grow with $n$: a single sensor
move changes the objective only through its own neighborhood. Both thresholds are
sharp for every instance, because the exact critical penalty attains them at the
extreme budgets $s\in\{0,n\}$. For the costs and edge-betweenness weights built
by the repository's \texttt{centrality} from nonnegative demands, every network
with $n\ge3$ satisfies $\bar\rho_{=}<2$, so $\rho=2$ is exact for every network
and every budget, while the rule $\rho=\sum_i|c_i|+1$ used in the notebook grows
from $3.8$ to $110$ across the networks in the repository. Finally, the exact per-instance critical penalty is
an explicit function of the coverage curve, and it vanishes at the budget the
notebook selects: there the penalty enforces nothing and only rescales the energy
landscape.
\end{abstract}

\section{Setup}
\label{sec:setup}

Let $G=(V,E)$ be a graph on $n:=|V|$ vertices, with a cost $c_i\in\mathbb R$ for
every vertex and a weight $w_{ij}\ge0$ for every undirected edge
$\{i,j\}\in E$. Write $N(i)$ for the neighbors of $i$,
\begin{equation}
W_i:=\sum_{j\in N(i)} w_{ij},
\qquad
a_i:=W_i-c_i .
\label{eq:Wa}
\end{equation}

\paragraph{Objective.}
For $\x\in\{0,1\}^n$, where $x_i=1$ places a sensor at $i$, set
\begin{equation}
f(\x):=\sum_{i\in V}c_ix_i+\sum_{\{i,j\}\in E}w_{ij}(1-x_i)(1-x_j).
\label{eq:obj}
\end{equation}
A pipe contributes its weight exactly when neither endpoint carries a sensor, so
$f$ trades the cost of the sensors against the weight of the pipes left
unmonitored. This is the objective of \texttt{MIQP}, \texttt{QUBO} and
\texttt{build\_Q\_matrix} in \texttt{src/formulations.py}. Write
$S(\x):=\{i:x_i=1\}$ and $k:=|S(\x)|$. For $i\in S$, let $R_i^-(\x)$ flip $x_i$
from $1$ to $0$; for $i\notin S$, let $R_i^+(\x)$ flip $x_i$ from $0$ to $1$.

\paragraph{Coverage curves.}
For $0\le k\le n$ define the equality and inequality optima
\begin{equation}
m_k:=\min\Bigl\{f(\x):\textstyle\sum_i x_i=k\Bigr\},
\qquad
M_k:=\min\Bigl\{f(\x):\textstyle\sum_i x_i\le k\Bigr\}=\min_{j\le k}m_j .
\label{eq:coverage}
\end{equation}
The sequence $(M_k)$ is the curve computed by \texttt{coverage} in
\texttt{src/experiments.py}, whose \texttt{MIQP} uses $\sum_ix_i\le s$.

\paragraph{Equality budget.}
For a budget $0\le s\le n$, the constrained problem has optimal value $m_s$ and
the penalized QUBO is
\begin{equation}
F_\rho(\x):=f(\x)+\rho\Bigl(\sum_{i\in V}x_i-s\Bigr)^2 .
\label{eq:Feq}
\end{equation}
This is $\x^\top Q\x+c_Q$ with $(Q,c_Q)$ returned by \texttt{build\_Q\_matrix}.

\paragraph{Inequality budget.}
For $s\ge0$, the constrained problem $\min\{f(\x):\sum_ix_i\le s\}$ has optimal
value $M_s$. Introduce $K$ slack bits $\y\in\{0,1\}^K$ and positive integer
coefficients $\beta_1,\dots,\beta_K$ such that
\begin{equation}
B(\y):=\sum_{\ell=1}^K\beta_\ell y_\ell,
\qquad
\{B(\y):\y\in\{0,1\}^K\}\supseteq\{0,1,\dots,s\}.
\label{eq:slack}
\end{equation}
The bounded binary encoding $K=\lfloor\log_2 s\rfloor+1$,
$\beta_\ell=2^{\ell-1}$ for $\ell<K$ and $\beta_K=s-2^{K-1}+1$ has range exactly
$\{0,\dots,s\}$; for $s=0$, $K=0$. The penalized QUBO is
\begin{equation}
G_\rho(\x,\y):=f(\x)+\rho\Bigl(\sum_{i\in V}x_i+B(\y)-s\Bigr)^2 .
\label{eq:Gle}
\end{equation}
Write $d(\x,\y):=\sum_ix_i+B(\y)-s$ for the residual of the slack equation.

\paragraph{Goal.}
We seek penalties guaranteeing \emph{strict exactness}: $F_\rho(\x)>m_s$ for
every $\x$ with $\sum_ix_i\neq s$, and $G_\rho(\x,\y)>M_s$ for every $(\x,\y)$
with $d(\x,\y)\neq0$. The per-instance critical penalties are
\begin{align}
\rseq(s)&:=\inf\{\rho>0:\text{every minimizer of }F_\rho\text{ is feasible}\},
\label{eq:rhostar-def}\\
\rsle(s)&:=\inf\{\rho>0:\text{every minimizer of }G_\rho\text{ has }d=0\}.
\nonumber
\end{align}

\begin{remark}[Spin units]
\label{rem:spin}
With $x_i=(z_i+1)/2$ one has $\sum_iz_i=2\sum_ix_i-n$, so
$\rho(\sum_ix_i-s)^2=\rho_z(\sum_iz_i-2s+n)^2$ with $\rho_z=\rho/4$. All
thresholds below are stated for the binary parameter $\rho$, as in the code; the
spin values are obtained by dividing by $4$, as in the balanced minimum-cut
analysis of \cite{RSDB2026}.
\end{remark}

\section{One-flip changes}
\label{sec:flips}

\begin{lemma}[One-flip changes]
\label{lem:flip}
For every $\x\in\{0,1\}^n$ and $i\in V$,
\begin{align}
f(R_i^+(\x))-f(\x)&=c_i-\sum_{j\in N(i):\,x_j=0}w_{ij}\in[\,c_i-W_i,\;c_i\,]
&&\text{if }x_i=0,
\label{eq:flip-add}\\
f(R_i^-(\x))-f(\x)&=-c_i+\sum_{j\in N(i):\,x_j=0}w_{ij}\in[\,-c_i,\;a_i\,]
&&\text{if }x_i=1.
\label{eq:flip-remove}
\end{align}
\end{lemma}

\begin{proof}
Only the terms of \eqref{eq:obj} that contain $x_i$ change: $c_ix_i$ and
$w_{ij}(1-x_i)(1-x_j)$ for $j\in N(i)$. The latter equals $w_{ij}$ when $x_i=x_j=0$
and $0$ otherwise. Placing a sensor at $i$ therefore adds $c_i$ and removes the
weight of every pipe from $i$ to an empty neighbor, while pipes to occupied
neighbors stay monitored; removing it reverses both effects. The intervals follow
from $0\le\sum_{j\in N(i):x_j=0}w_{ij}\le W_i$.
\end{proof}

The change in Lemma~\ref{lem:flip} depends only on the neighborhood of $i$. This is
the structural difference from the balanced minimum-cut objective of
\cite{RSDB2026}, where removing a vertex from a set of size $k$ changes the cut by
up to $2(k-1)$, which grows with $n$ and forces the threshold $\rho>n$.

\begin{lemma}[Incident edge betweenness]
\label{lem:betweenness}
Let $w_{ij}$ be the normalized edge betweenness,
\[
w_{ij}=\frac{2}{n(n-1)}\sum_{\{u,v\}\subseteq V}\frac{\sigma_{uv}(ij)}{\sigma_{uv}},
\]
where $\sigma_{uv}$ counts the shortest $u$--$v$ paths, $\sigma_{uv}(ij)$ those
using $\{i,j\}$, and pairs in different components contribute $0$. This is
\texttt{networkx.edge\_betweenness\_centrality} as called by \texttt{centrality}.
Then $W_i\le 2(n-1)/n<2$ for every $i$, with equality at the center of a star.
\end{lemma}

\begin{proof}
Summing over the edges at $i$,
$W_i=\frac{2}{n(n-1)}\sum_{\{u,v\}}\sigma_{uv}^{-1}\sum_P|P\cap\delta(i)|$, where
$P$ ranges over the shortest $u$--$v$ paths and $\delta(i)$ is the set of edges at
$i$. A shortest path is simple, so it contains at most one edge at $i$ when $i$ is
an endpoint and at most two otherwise. There are $n-1$ pairs containing $i$ and
$\binom{n-1}{2}$ pairs not containing it, so
\[
W_i\le\frac{2}{n(n-1)}\Bigl[(n-1)+2\binom{n-1}{2}\Bigr]
=\frac{2}{n(n-1)}(n-1)(n-1)=\frac{2(n-1)}{n}.
\]
At the center of a star every pair containing the center uses exactly one edge at
it and every pair of leaves passes through it using exactly two, so both bounds are
attained.
\end{proof}

\section{Equality budget}
\label{sec:equality}

\begin{theorem}[Strict exactness for $\rho>\bar\rho_{=}$]
\label{thm:eq}
Let $0\le s\le n$ and
\begin{equation}
\rbeq:=\max_{i\in V}\max\{c_i,\;W_i-c_i\}.
\label{eq:rbeq}
\end{equation}
If $\rho>\rbeq$, then
$\argmin_{\x\in\{0,1\}^n}F_\rho=\argmin\{f(\x):\sum_ix_i=s\}$; in particular every
minimizer of $F_\rho$ is feasible and $F_\rho(\x)>m_s$ for every infeasible $\x$.
Note that $\rbeq\ge\max_iW_i/2\ge0$.
\end{theorem}

\begin{proof}
On feasible points the penalty vanishes, so $\min_\x F_\rho\le m_s$. Suppose
$\tilde\x$ minimizes $F_\rho$ and is infeasible, with $\tilde S:=S(\tilde\x)$ and
$k=|\tilde S|\neq s$.

If $k=s+r$, $r\ge1$, pick $i\in\tilde S$ and set $\x'=R_i^-(\tilde\x)$. By
\eqref{eq:flip-remove}, $f(\x')-f(\tilde\x)\le a_i\le\rbeq$. The penalty changes by
$\rho\bigl[(r-1)^2-r^2\bigr]=-\rho(2r-1)$, hence
\begin{equation}
F_\rho(\x')-F_\rho(\tilde\x)
\le a_i-\rho(2r-1)
<\rbeq-\rbeq(2r-1)
=-2\rbeq(r-1)\le0,
\label{eq:eq-drop}
\end{equation}
where the strict inequality uses $\rho>\rbeq$ and $2r-1\ge1$. This contradicts the
optimality of $\tilde\x$. If $k=s-r$, $r\ge1$, then $k<n$, so some $i\notin\tilde S$
exists; set $\x'=R_i^+(\tilde\x)$. By \eqref{eq:flip-add},
$f(\x')-f(\tilde\x)\le c_i\le\rbeq$, the penalty drop is again $\rho(2r-1)$, and
\eqref{eq:eq-drop} holds with $c_i$ in place of $a_i$.

Thus every minimizer is feasible; on the feasible set $F_\rho=f$, so the
minimizers of $F_\rho$ are exactly those of the constrained problem, and any
infeasible $\x$ satisfies $F_\rho(\x)>\min F_\rho=m_s$.
\end{proof}

\begin{theorem}[Boundary case $\rho=\bar\rho_{=}$]
\label{thm:eq-boundary}
Let $\rho=\rbeq$ and $\tilde\x\in\argmin F_{\rbeq}$. A finite sequence of one-bit
repairs ($R_i^-$ while $\sum_ix_i>s$, $R_i^+$ while $\sum_ix_i<s$) produces a
feasible $\x^0$ with $F_{\rbeq}(\x^0)\le F_{\rbeq}(\tilde\x)$ and $f(\x^0)=m_s$.
Hence $\min F_{\rbeq}=m_s$, although infeasible optimal ties may occur.
\end{theorem}

\begin{proof}
For any $\x$ with $\sum_ix_i=s+r$, $r\ge1$, and $i\in S(\x)$, \eqref{eq:eq-drop}
with $\rho=\rbeq$ gives $F_{\rbeq}(R_i^-(\x))-F_{\rbeq}(\x)\le-2\rbeq(r-1)\le0$;
the case $\sum_ix_i=s-r$ is analogous. Each step is nonworsening and reduces
$|\sum_ix_i-s|$ by one, so a feasible $\x^0$ with
$F_{\rbeq}(\x^0)\le F_{\rbeq}(\tilde\x)\le m_s$ is reached, and feasibility forces
$f(\x^0)=F_{\rbeq}(\x^0)\ge m_s$.
\end{proof}

The proof of Theorem~\ref{thm:eq} only needs \emph{one} repairable vertex, which
yields a budget-dependent refinement. Write $a_{[1]}\ge\dots\ge a_{[n]}$ and
$c_{[1]}\ge\dots\ge c_{[n]}$ for the sorted values, and set
$a_{[n+1]}=c_{[n+1]}:=-\infty$.

\begin{corollary}[Budget-dependent threshold]
\label{cor:eq-budget}
Let
\begin{equation}
\rbeq(s):=\max\bigl\{0,\;a_{[s+1]},\;c_{[n-s+1]}\bigr\}\le\rbeq .
\label{eq:rbeq-s}
\end{equation}
If $\rho>\rbeq(s)$, the conclusions of Theorem~\ref{thm:eq} hold.
\end{corollary}

\begin{proof}
If $k=s+r$, the set $\tilde S$ has at least $s+1$ elements, so it contains some
$i$ with $a_i\le a_{[s+1]}$; repairing that $i$ gives
$F_\rho(\x')-F_\rho(\tilde\x)\le a_{[s+1]}-\rho(2r-1)\le a_{[s+1]}-\rho<0$. If
$k=s-r$, the complement $V\setminus\tilde S$ has at least $n-s+1$ elements, so it
contains some $i$ with $c_i\le c_{[n-s+1]}$, and the same estimate holds. The
inequality $\rbeq(s)\le\rbeq$ is immediate, since $a_{[s+1]}\le\max_ia_i$,
$c_{[n-s+1]}\le\max_ic_i$ and $\rbeq\ge0$.
\end{proof}

\begin{proposition}[Sharpness of the equality threshold]
\label{prop:eq-sharp}
\leavevmode
\begin{enumerate}
\item[(a)] \emph{Every instance attains $\rbeq$.} For every instance,
$\rseq(0)=\max\{0,\max_ia_i\}$ and $\rseq(n)=\max\{0,\max_ic_i\}$. Consequently
\[
\rbeq=\max_{0\le s\le n}\rseq(s),
\]
so no budget-independent penalty below $\rbeq$ is exact for every budget of a given
network, and Corollary~\ref{cor:eq-budget} is attained at $s\in\{0,n\}$.
\item[(b)] \emph{Pipe term, every budget.} For $s\ge0$ and $w>0$, let $G$ consist
of $s+1$ disjoint pipes of weight $w$ with $c\equiv0$. Then $\rbeq=w$, and for
every $\rho<w$ every minimizer of $F_\rho$ is infeasible; at $\rho=w$ infeasible
optimal ties occur.
\item[(c)] \emph{Cost term, every budget.} For $s\ge1$ and $c>0$, let $G$ consist
of $n=s+1$ isolated vertices of cost $c$. Then $\rbeq(s)=\rbeq=c$, and for every
$\rho<c$ every minimizer of $F_\rho$ is infeasible.
\end{enumerate}
\end{proposition}

\begin{proof}
(a) For $s=0$ the only feasible point is $\x=\mathbf 0$, with
$m_0=\sum_{\{i,j\}\in E}w_{ij}$. Placing one sensor at $i$ gives, by
\eqref{eq:flip-add} with every neighbor empty, $F_\rho=m_0-a_i+\rho$, which beats
$m_0$ whenever $\rho<a_i$; hence $\rseq(0)\ge\max\{0,\max_ia_i\}$, and
Corollary~\ref{cor:eq-budget} at $s=0$ gives the matching upper bound. For $s=n$
the only feasible point is $\ones$, with $m_n=\sum_ic_i$; removing the sensor at
$i$ leaves every pipe monitored, so $F_\rho=m_n-c_i+\rho$, and the same two
arguments give $\rseq(n)=\max\{0,\max_ic_i\}$. Finally
$\rbeq=\max\{\rseq(0),\rseq(n)\}\le\max_s\rseq(s)\le\rbeq$, the last inequality by
Theorem~\ref{thm:eq}.

(b) Here $n=2s+2$, $W_i=w$ and $c_i=0$, so $\rbeq=w$. A feasible point places $s$
sensors and monitors at most $s$ pipes, so $m_s=w$, attained with one sensor on
each of $s$ pipes. The point $\x^\bullet$ with one sensor on every pipe has
$f(\x^\bullet)=0$ and $F_\rho(\x^\bullet)=\rho<w=m_s$, so $\x^\bullet$ beats every
feasible point and all minimizers are infeasible. At $\rho=w$,
$F_w(\x^\bullet)=m_s$ is an infeasible tie.

(c) There are no pipes, so $a_i=-c$, $c_{[n-s+1]}=c_{[2]}=c$ and
$\rbeq(s)=\rbeq=c$. Feasible points cost $m_s=sc$, while $s-1$ sensors give
$F_\rho=(s-1)c+\rho<sc$ for $\rho<c$.
\end{proof}

\section{Inequality budget}
\label{sec:inequality}

\begin{theorem}[Strict exactness for $\rho>\bar\rho_{\le}$]
\label{thm:le}
Let $s\ge0$, let $B$ satisfy \eqref{eq:slack}, and let
\begin{equation}
\rble:=\max\Bigl\{0,\;\max_{i\in V}(W_i-c_i)\Bigr\}.
\label{eq:rble}
\end{equation}
If $\rho>\rble$, then every minimizer $(\tilde\x,\tilde\y)$ of $G_\rho$ satisfies
$d(\tilde\x,\tilde\y)=0$, hence $\sum_i\tilde x_i\le s$; moreover
$\min G_\rho=M_s$, and the $\x$-parts of the minimizers of $G_\rho$ are exactly the
minimizers of $\{f(\x):\sum_ix_i\le s\}$.
\end{theorem}

\begin{proof}
For every $\x$ with $\sum_ix_i\le s$, \eqref{eq:slack} provides $\y_\x$ with
$B(\y_\x)=s-\sum_ix_i$, so $G_\rho(\x,\y_\x)=f(\x)$ and $\min G_\rho\le M_s$.
Suppose $(\tilde\x,\tilde\y)$ minimizes $G_\rho$ with $\tilde d:=d(\tilde\x,\tilde\y)\neq0$.

If $\sum_i\tilde x_i\le s$, then
$G_\rho(\tilde\x,\y_{\tilde\x})=f(\tilde\x)<f(\tilde\x)+\rho\tilde d^2=G_\rho(\tilde\x,\tilde\y)$
because $\rho>0$, a contradiction.

If $\sum_i\tilde x_i=s+r$ with $r\ge1$, then $\tilde d\ge r\ge1$ because
$B\ge0$. Pick $i\in S(\tilde\x)$, set $\x'=R_i^-(\tilde\x)$ and keep $\tilde\y$. By
\eqref{eq:flip-remove}, $f(\x')-f(\tilde\x)\le a_i\le\rble$, and the penalty
changes by $\rho\bigl[(\tilde d-1)^2-\tilde d^2\bigr]=-\rho(2\tilde d-1)$, so
\begin{equation}
G_\rho(\x',\tilde\y)-G_\rho(\tilde\x,\tilde\y)\le a_i-\rho(2\tilde d-1)
<\rble-\rble(2\tilde d-1)\le0,
\label{eq:le-drop}
\end{equation}
again a contradiction. Hence every minimizer has $d=0$ and thus
$\sum_i\tilde x_i\le s$ and $G_\rho(\tilde\x,\tilde\y)=f(\tilde\x)\ge M_s$; combined
with $\min G_\rho\le M_s$ this gives $\min G_\rho=M_s$ and the optimality of
$\tilde\x$. Conversely every optimal $\x$ extends to the minimizer $(\x,\y_\x)$.
\end{proof}

\begin{theorem}[Boundary case $\rho=\bar\rho_{\le}$]
\label{thm:le-boundary}
If $\rble>0$, then $\min G_{\rble}=M_s$, although infeasible optimal ties may occur.
\end{theorem}

\begin{proof}
Let $(\tilde\x,\tilde\y)$ be a minimizer. While $\sum_ix_i>s$ we have $d\ge1$, and
\eqref{eq:le-drop} with $\rho=\rble$ shows that removing any sensor does not
increase $G_{\rble}$. After finitely many removals $\sum_ix_i\le s$, and replacing
$\y$ by $\y_\x$ gives a point with value $f(\x)\ge M_s$ and no larger than
$G_{\rble}(\tilde\x,\tilde\y)\le M_s$.
\end{proof}

\begin{corollary}[Budget-dependent threshold]
\label{cor:le-budget}
If $\rho>\rble(s):=\max\{0,a_{[s+1]}\}\le\rble$, the conclusions of
Theorem~\ref{thm:le} hold.
\end{corollary}

\begin{proof}
In the case $\sum_i\tilde x_i=s+r$, the set $S(\tilde\x)$ has at least $s+1$
elements and contains some $i$ with $a_i\le a_{[s+1]}$; repair that vertex in
\eqref{eq:le-drop}.
\end{proof}

\begin{proposition}[Sharpness of the inequality threshold]
\label{prop:le-sharp}
(a) For every instance, $\rsle(0)=\rble$, so $\rble=\max_{s\ge0}\rsle(s)$.
(b) For the $s+1$ disjoint pipes of Proposition~\ref{prop:eq-sharp}(b),
$\rble=w$, $M_s=w$, and for every $\rho<w$ every minimizer of $G_\rho$ places
more than $s$ sensors.
\end{proposition}

\begin{proof}
(a) For $s=0$ the encoding is empty ($K=0$), so $G_\rho=F_\rho$ and
$\rsle(0)=\rseq(0)=\max\{0,\max_ia_i\}=\rble$ by
Proposition~\ref{prop:eq-sharp}(a); Theorem~\ref{thm:le} gives $\rsle(s)\le\rble$
for every $s$.
(b) A point with at most $s$ sensors monitors at most $s$ pipes, so $M_s=w$. The
point $\x^\bullet$ with one sensor per pipe has $\min_\y G_\rho(\x^\bullet,\y)=\rho$,
attained at $\y=\mathbf0$, which is below $M_s$ when $\rho<w$; hence no point
with at most $s$ sensors is a minimizer.
\end{proof}

\begin{remark}[Why the cost term disappears]
\label{rem:onesided}
Under the inequality budget, a deficit $\sum_ix_i<s$ is absorbed by the slack at no
cost, so only the removal direction \eqref{eq:flip-remove} must be repaired, and
$c_i$ no longer enters the threshold. Equivalently, for every $\x$ and $\rho>0$,
$\min_\y G_\rho(\x,\y)=f(\x)+\rho\max\{0,\sum_ix_i-s\}^2$: the slack turns the
two-sided penalty into a one-sided one. The price is $K=\lfloor\log_2s\rfloor+1$
extra variables, coupled to every $x_i$ with strength $\rho\beta_\ell$, so the QUBO
remains a complete graph on $n+K$ variables.
\end{remark}

\section{Water distribution networks}
\label{sec:wdn}

The repository builds $c$ and $w$ with \texttt{centrality} in
\texttt{src/network.py}: $c_i=q_i/\max_jq_j+\deg(i)/(n-1)^2$, where $q_i\ge0$ is
the water demand at $i$ and $\deg(i)/(n-1)$ is \texttt{networkx}'s degree
centrality divided by $n-1$ once more, and $w_{ij}$ is the normalized edge
betweenness of Lemma~\ref{lem:betweenness}.

\begin{corollary}[A size-independent exact penalty]
\label{cor:wdn}
Assume $q_i\ge0$ for all $i$ and $\max_jq_j>0$. Then $0\le c_i\le1+\frac1{n-1}$ and
$W_i\le\frac{2(n-1)}{n}$, so
\[
\rbeq\le\max\Bigl\{1+\tfrac1{n-1},\;\tfrac{2(n-1)}{n}\Bigr\},
\qquad
\rble\le\tfrac{2(n-1)}{n}.
\]
In particular, for $n\ge3$ the penalty $\rho=2$ is strictly exact for every
network, every budget $s$, and both the equality and the inequality budget. The
notebook's rule $\rho_{\rm nb}:=\sum_i|c_i|+1$ satisfies $\rho_{\rm nb}\ge2$ and is
therefore also exact for $n\ge3$, but
$\rho_{\rm nb}\ge1+\sum_iq_i/\max_jq_j$ grows linearly in $n$ for comparable
demands.
\end{corollary}

\begin{proof}
The demand term lies in $[0,1]$ and the degree term in $[0,\frac1{n-1}]$, since
$\deg(i)\le n-1$. Lemma~\ref{lem:betweenness} bounds $W_i$, and
$W_i-c_i\le W_i$. For $n\ge3$, $1+\frac1{n-1}\le\frac32<2$ and
$\frac{2(n-1)}n<2$, so Theorems~\ref{thm:eq} and~\ref{thm:le} apply at $\rho=2$.
If $i^\ast$ has the largest demand then $c_{i^\ast}\ge1$, so
$\rho_{\rm nb}\ge c_{i^\ast}+1\ge2$.
\end{proof}

\begin{remark}[Degree term]
\label{rem:degree}
If the second division by $n-1$ is removed, so that the degree term is
\texttt{networkx}'s degree centrality in $[0,1]$ (SUSPECT-1 in
\texttt{docs/code\_review.md}), the same argument gives $c_i\le2$ and $\rbeq\le2$,
so every $\rho>2$ is exact.
\end{remark}

Table~\ref{tab:networks} evaluates the thresholds on every network in
\texttt{Data/WDN\_Data}. In every network $\rbeq=\max_ic_i\approx1$, set by the
vertex with the largest demand, regardless of size, while $\rho_{\rm nb}$ is $4$ to
$110$ times larger.

\begin{table}[h]
\centering
\begin{tabular}{lrrcccccc}
\toprule
Network & $n$ & $|E|$ & $\max_ic_i$ & $\max_iW_i$ & $\frac{2(n-1)}{n}$ & $\rho_{\rm nb}$ & $\rbeq$ & $\rble$\\
\midrule
Test     & 18  & 18  & 1.0069 & 0.2353 & 1.8889 & 3.821   & 1.0069 & 0.1636\\
Apulia   & 24  & 34  & 1.0057 & 0.5709 & 1.9167 & 4.473   & 1.0057 & 0.4493\\
Fossolo  & 37  & 58  & 1.0023 & 0.4205 & 1.9459 & 5.581   & 1.0023 & 0.3386\\
ZJ       & 114 & 162 & 1.0002 & 1.0506 & 1.9825 & 23.659  & 1.0002 & 0.9130\\
Modena   & 272 & 317 & 1.0000 & 0.6186 & 1.9926 & 43.980  & 1.0000 & 0.4311\\
Kentucky & 856 & 982 & 1.0000 & 0.5498 & 1.9977 & 109.987 & 1.0000 & 0.4944\\
\bottomrule
\end{tabular}
\caption{Penalty thresholds for the networks in \texttt{Data/WDN\_Data}, with costs
and weights from \texttt{centrality}. $\rho_{\rm nb}=\sum_i|c_i|+1$ is the value
used in the notebook and stored in \texttt{logs/}. Values are rounded to four
decimals; $\max_ic_i$ exceeds $1$ only through the degree term, which is at most
$1/(n-1)$.}
\label{tab:networks}
\end{table}

\section{Per-instance critical penalty}
\label{sec:rhostar}

Because the configuration space is finite, the critical penalties are explicit
functions of the coverage curves \eqref{eq:coverage}, in the same way as the
identity for $\rho^\star$ in the balanced minimum-cut analysis of
\cite{RSDB2026}.

\begin{proposition}[Critical penalty from the coverage curve]
\label{prop:rhostar}
For every instance,
\begin{align}
\rseq(s)&=\max\Bigl\{0,\;\max_{k\neq s}\frac{m_s-m_k}{(k-s)^2}\Bigr\},
\label{eq:rhostar-eq}\\
\rsle(s)&=\max\Bigl\{0,\;\max_{k>s}\frac{M_s-M_k}{(k-s)^2}\Bigr\}.
\label{eq:rhostar-le}
\end{align}
\end{proposition}

\begin{proof}
For \eqref{eq:rhostar-eq}: every minimizer of $F_\rho$ is feasible exactly when
$f(\x)+\rho(\sum_ix_i-s)^2>m_s$ for every infeasible $\x$. For a fixed count
$k\neq s$ the most demanding $\x$ is one with $f(\x)=m_k$, so the condition reads
$\rho>(m_s-m_k)/(k-s)^2$ for all $k\neq s$; at equality an infeasible tie exists.

For \eqref{eq:rhostar-le}: if $\sum_ix_i\le s$ then $\min_\y G_\rho(\x,\y)=f(\x)$,
attained only where $d=0$ when $\rho>0$. If $\sum_ix_i=k>s$ then $d\ge k-s$ with
equality at $\y=\mathbf0$, so $\min_\y G_\rho(\x,\y)=f(\x)+\rho(k-s)^2$. Hence every
minimizer has $d=0$ if and only if $\rho>0$ and $\rho>(M_s-m_k)/(k-s)^2$ for every
$k>s$. Replacing $m_k$ by $M_k\le m_k$ cannot lower any term; conversely, if
$M_k=m_j$ with $s<j\le k$ then $(M_s-M_k)/(k-s)^2\le(M_s-m_j)/(j-s)^2$, and if
$M_k=M_s$ the term is $0$. The two maxima therefore coincide.
\end{proof}

The notebook chooses the budget as the first sensor count at which the coverage
curve reaches its minimum (cell 14). At that budget no penalty is needed.

\begin{proposition}[Budget at the minimum of the coverage curve]
\label{prop:sstar}
Let $s^\ast$ be the smallest minimizer of $k\mapsto M_k$. Then
$m_{s^\ast}=M_{s^\ast}=\min_km_k$, and $\rseq(s^\ast)=\rsle(s^\ast)=0$: every
$\rho>0$ is strictly exact at $s^\ast$, for both budgets.
\end{proposition}

\begin{proof}
The sequence $(M_k)$ is nonincreasing with $M_n=\min_km_k$. If
$m_{s^\ast}>M_{s^\ast}$ then $s^\ast\ge1$ and
$M_{s^\ast-1}=M_{s^\ast}$, contradicting the minimality of $s^\ast$. Hence
$m_{s^\ast}=\min_km_k$, every term $m_{s^\ast}-m_k$ in \eqref{eq:rhostar-eq} is
nonpositive, and $M_{s^\ast}-M_k=0$ for all $k>s^\ast$ in \eqref{eq:rhostar-le}.
\end{proof}

At $s^\ast$ the unconstrained minimum of $f$ already places $s^\ast$ sensors, and
the penalty only breaks ties between budgets. Its magnitude then plays no role in
exactness and only rescales the landscape seen by simulated annealing, tabu search
and the annealer.

\paragraph{Exact values, numerically.}
For the two networks small enough to enumerate, Table~\ref{tab:rhostar} reports
the exact critical penalties over all budgets, computed from \eqref{eq:rhostar-eq}
and \eqref{eq:rhostar-le} with the coverage curves obtained by evaluating all
$2^n$ placements. For Apulia, $s^\ast=11$ is the budget stored in
\texttt{logs/Apulia}, and $m_{11}=1.300182$ equals, to ten digits, the best tabu
search energy in \texttt{Apulia\_TabuSearch.json}.

\begin{table}[h]
\centering
\begin{tabular}{lrrcccccc}
\toprule
Network & $n$ & $s^\ast$ & $m_{s^\ast}$ & $\rseq(s^\ast)$ & $\max_{0<s<n}\rseq(s)$ & $\operatorname{med}_s\rseq(s)$ & $\rseq(n)=\rbeq$ & $\rsle(0)=\rble$\\
\midrule
Test   & 18 & 6  & 0.8115 & 0 & 0.1418 & 0.1210 & 1.0069 & 0.1636\\
Apulia & 24 & 11 & 1.3002 & 0 & 0.3546 & 0.1109 & 1.0057 & 0.4493\\
\bottomrule
\end{tabular}
\caption{Exact critical penalties by enumeration of all $2^n$ placements. The
maximum over intermediate budgets is attained at $s=1$ in both networks, and
$\rsle(s)=0$ for every $s\ge s^\ast$.}
\label{tab:rhostar}
\end{table}

\begin{remark}[Computing $\rho^\star$ at scale]
\label{rem:compute}
Proposition~\ref{prop:rhostar} reduces $\rho^\star$ to $n+1$ constrained optima.
For $n\lesssim25$ they follow from exhaustive enumeration, as above. For larger
networks they are $n+1$ MIQP solves with $\sum_ix_i=k$; for the inequality budget
the curve $(M_k)$ returned by \texttt{coverage} already suffices. Since
\texttt{coverage} uses $\sum_ix_i\le s$ (BUG-03 in
\texttt{docs/code\_review.md}), it returns $M$, not $m$, so \eqref{eq:rhostar-eq}
requires the equality version.
\end{remark}

\section{Synthesis and open problems}
\label{sec:synthesis}

Three penalties describe the same QUBO:
\begin{itemize}
\item \textbf{Instance threshold.} $\rho>\rbeq$ (equality) or $\rho>\rble$
(inequality) is exact for every budget (Theorems~\ref{thm:eq}--\ref{thm:le-boundary}),
and no smaller budget-independent penalty is (Propositions~\ref{prop:eq-sharp}
and~\ref{prop:le-sharp}). The thresholds depend only on the neighborhood of each
vertex, and for the repository's networks they are below $2$ independently of $n$
(Corollary~\ref{cor:wdn}).
\item \textbf{Per-instance critical penalty.} $\rho^\star(s)$ is an explicit
function of the coverage curve (Proposition~\ref{prop:rhostar}); it vanishes at
the budget the notebook selects (Proposition~\ref{prop:sstar}) and stays below
$0.36$ at every intermediate budget of the enumerated networks.
\item \textbf{Notebook rule.} $\rho_{\rm nb}=\sum_i|c_i|+1$ is exact but grows
linearly in $n$, reaching $110$ on Kentucky.
\end{itemize}
There is no conflict between these values: $\rho_{\rm nb}$ is exact because it
exceeds $\rbeq$, and $\rbeq$ is the worst case of $\rho^\star(s)$ over budgets.

\paragraph{Practical guidance.} Use $\rho=2$ for every network with $n\ge3$, or
$\rbeq(s)(1+\varepsilon)$ from Corollary~\ref{cor:eq-budget} when a smaller value is
wanted for a fixed budget. At $s^\ast$, any small $\rho>0$ is exact. The penalty
sets the scale of $Q$: the diagonal entries are
$c_i-W_i+\rho(1-2s)$, which for Modena ($s=87$) is about $-7608$ at
$\rho_{\rm nb}=43.98$, the value printed in the notebook, against about $-346$ at
$\rho=2$, a dynamic range $22$ times smaller. Whether this improves the samplers is
an empirical question that the logs cannot answer, since every run used
$\rho_{\rm nb}$; it is most pressing for quantum annealing, which found no feasible
sample on ZJ.

\paragraph{Open problems.} First, the effect of $\rho$ on heuristic and quantum
samplers, as opposed to exactness, which is settled here: a smaller penalty
flattens the barriers between budgets, which may help or hurt mixing. Second,
encodings of the inequality budget without slack bits, such as one-sided penalties
that are not quadratic, are not covered by Theorem~\ref{thm:le}. Third, problems
with several budgets, for example per-district sensor limits, add one penalty per
constraint; the one-flip argument extends when every vertex appears in a bounded
number of constraints, but the interaction between penalties is not analyzed here.

\paragraph{Relation to prior work.}
The one-bit repair arguments follow the QUBO exactness pattern of Quintero, Bernal,
Terlaky, and Zuluaga \cite{QBTZ2021}, and the structure of this note follows the
balanced minimum-cut analysis of \cite{RSDB2026}. Slack-variable encodings of
inequality constraints in QUBOs are standard; see \cite{Lucas2014}.

\paragraph{Reproducibility.}
Every statement above is checked numerically by
\nolinkurl{docs/penalty_parameter/verify_penalty_thresholds.py}. Run it from the
repository root. It checks:
\begin{itemize}
\item Lemmas~\ref{lem:flip}--\ref{lem:betweenness} on random graphs;
\item Theorems~\ref{thm:eq}--\ref{thm:le-boundary},
Corollaries~\ref{cor:eq-budget} and~\ref{cor:le-budget},
Propositions~\ref{prop:eq-sharp}(a), \ref{prop:le-sharp}(a), \ref{prop:rhostar}
and~\ref{prop:sstar}, by exhaustive enumeration of the penalized QUBOs, including
the slack bits, on $300$ random connected graphs with $n\le9$ (seed
$20261007$), half scored with \texttt{centrality} and half with generic
$c_i\sim U[-1,2]$, $w_{ij}\sim U[0,2]$;
\item the sharpness families of Propositions~\ref{prop:eq-sharp}(b)--(c)
and~\ref{prop:le-sharp}(b).
\end{itemize}
It also computes Tables~\ref{tab:networks} and~\ref{tab:rhostar}, which it writes
to \texttt{penalty\_thresholds\_results.json}.

\begin{thebibliography}{3}
\bibitem{Lucas2014}
A.~Lucas, ``Ising formulations of many NP problems,'' \emph{Frontiers in Physics},
2 (2014), p.~5.
\url{https://doi.org/10.3389/fphy.2014.00005}.

\bibitem{QBTZ2021}
R.~Quintero, D.~Bernal, T.~Terlaky, and L.~F.~Zuluaga,
``Characterization of QUBO reformulations for the maximum $k$-colorable subgraph
problem,'' \emph{arXiv preprint} arXiv:2101.09462, 2021.
\url{https://arxiv.org/abs/2101.09462}.

\bibitem{RSDB2026}
A.~Ramesh, B.~Sundar, M.~Dupont, and D.~E.~Bernal Neira,
``Quantum preconditioning for constrained optimization problems,''
\emph{arXiv preprint} arXiv:2608.28842, 2026, Appendix~A.
\url{https://arxiv.org/abs/2608.28842}.
\end{thebibliography}

\end{document}
